% Start preamble
\documentclass[12pt,a4paper]{article}
\usepackage{geometry}
 \geometry{
 a4paper,
 total={170mm,257mm},
 left=20mm,
 top=20mm,
 }
\usepackage{currfile}
\usepackage[utf8]{inputenc}
\usepackage[T1]{fontenc}
\usepackage[pdftex]{graphicx}
\graphicspath{{./}}
\usepackage{enumitem}
\usepackage{pdfpages}
\usepackage{hyperref}
\usepackage{tikz}
\usepackage{attachfile}
\usepackage{epstopdf}
\usepackage{array}
\usepackage{multirow}
\usepackage{multicol}
\usepackage{float}
%\usepackage[table]{xcolor,colorbl}
\setlength{\textwidth}{16cm}
\setlength{\oddsidemargin}{-0.5cm}
\setlength{\evensidemargin}{-0.5cm}
%\setlenght{\headsep}{0cm}
\setlength\parindent{0pt}
%\setlength{\extrarowheight}{3pt}
\usepackage{listings}
%\usepackage{xcolor}
\input{arduinoLanguage.tex}
%%%%%% Counting oppgaves %%%%%%
 \newcount\questnum \questnum=0
 \def\oppgave{
            \advance\questnum by 1
            \ifnum \questnum > 0
                 \hrule
                 \vskip 5pt
                 \leftline{Oppgave \the\questnum}
                 \vskip 5pt \fi
		}
% End preamble
\begin{document}
\title{Oppgaer til Turtall og Posisjon}
\author{Faglærer: Fred-Olav Mosdal 90507684\\}
\maketitle
\oppgave{}%1
\vskip 2.5pt 
Hva er forskjellen mellom en turtallsgiver og en posisjonsgiver?
\vskip 5pt 
\begin{tikzpicture}
	\draw[step=0.5cm,gray!20,very thin]  grid (17,4) ;
\end{tikzpicture}
\vskip 2.5pt 
\oppgave{}%2
\vskip 2.5pt 
Hvilken type posisjonsgiver brukes vanligvis til å måle rotasjonsposisjonen til en motor?
\vskip 5pt 
\begin{tikzpicture}
	\draw[step=0.5cm,gray!20,very thin]  grid (17,10) ;
\end{tikzpicture}
\vskip 2.5pt 
\oppgave{}%3
\vskip 2.5pt 
Hvordan kan en turtallsgiver brukes til å overvåke hastigheten til en motor?
\vskip 5pt 
\begin{tikzpicture}
	\draw[step=0.5cm,gray!20,very thin]  grid (17,10) ;
\end{tikzpicture}
\vskip 2.5pt 
\oppgave{}%4
\vskip 2.5pt 
Hva er en enkoder, og hva brukes den til?
\vskip 5pt 
\begin{tikzpicture}
	\draw[step=0.5cm,gray!20,very thin]  grid (17,6) ;
\end{tikzpicture}
\vskip 2.5pt 
\oppgave{}%5
\vskip 2.5pt 
Hva er forskjellen mellom en absolutt og en inkrementell enkoder?


\vskip 5pt 
\begin{tikzpicture}
	\draw[step=0.5cm,gray!20,very thin]  grid (17,4) ;
\end{tikzpicture}
\vskip 2.5pt 
\oppgave{}%6
\vskip 2.5pt 
Søk opp manualen til Omron E6B2-CWZ6C-1000 og svar på følgende:

\begin{itemize}
	\item Hva er oppløsningen på denne enkoderen, og hva betyr det?

\begin{tikzpicture}
	\draw[step=0.5cm,gray!20,very thin]  grid (16,3) ;
\end{tikzpicture}

	\item Hvilken type utgang har denne enkoderen?

\begin{tikzpicture}
	\draw[step=0.5cm,gray!20,very thin]  grid (16,3) ;
\end{tikzpicture}

	\item Tegn en skisse for hvordan den kobles til en PLS

\begin{tikzpicture}
	\draw[step=0.5cm,gray!20,very thin]  grid (16,3) ;
\end{tikzpicture}

	\item Hvordan kan enkoderen brukes til å måle rotasjonshastigheten til en aksel?

\begin{tikzpicture}
	\draw[step=0.5cm,gray!20,very thin]  grid (16,3) ;
\end{tikzpicture}

\end{itemize}

\vskip 5pt 
\begin{tikzpicture}
	\draw[step=0.5cm,gray!20,very thin]  grid (17,6) ;
\end{tikzpicture}
\vskip 2.5pt 
\oppgave{}%7
\vskip 2.5pt 
Hva brukes en lineærgiver til?
\vskip 5pt 
\begin{tikzpicture}
	\draw[step=0.5cm,gray!20,very thin]  grid (17,6) ;
\end{tikzpicture}
\vskip 2.5pt 
\oppgave{}%7
\vskip 2.5pt 
Hva er en linelengdegiver?
%$$\includegraphics[width=0.4\textwidth]{../output/nogpl/trukkitank.png}$$
\vskip 5pt 
\begin{tikzpicture}
	\draw[step=0.5cm,gray!20,very thin]  grid (17,6) ;
\end{tikzpicture}
\vskip 2.5pt 
\oppgave{}%8
\vskip 2.5pt 
En RLT brukes til å måle lengden av en vaier på en heis. Trommelen roterer 50 ganger og hver rotasjon tilsvarer en lengde på 0,2 meter. Hva er lengden på vaieren?
%$$\includegraphics[width=0.4\textwidth]{../output/nogpl/trukkitank.png}$$
\vskip 5pt 
\begin{tikzpicture}
	\draw[step=0.5cm,gray!20,very thin]  grid (17,6) ;
\end{tikzpicture}
\vskip 2.5pt 
\oppgave{}%9
\vskip 2.5pt 
En RLT brukes til å overvåke spenningen i en vaier på en kran. Vaieren skal ha en spenning på 1000 Newton. Hvis RLT-en registrerer en spenning på 1200 Newton, hva er prosentvis økning i spenningen?
%$$\includegraphics[width=0.4\textwidth]{../output/nogpl/trukkitank.png}$$
\vskip 5pt 
\begin{tikzpicture}
	\draw[step=0.5cm,gray!20,very thin]  grid (17,6) ;
\end{tikzpicture}
\vskip 2.5pt 
\end{document}
